
\documentclass{article}



\usepackage{graphicx,comment,framed}
\usepackage{roundbox}
\usepackage{fancybox}
\usepackage{tikz}
\usepackage{color}
\usepackage[hidelinks]{hyperref}
\usepackage{framed}
\usepackage{amsthm,amssymb,amsmath}
%\usepackage[colorlinks,linkcolor=blue,citecolor=blue]{hyperref}
\definecolor{shadecolor}{cmyk}{0,0,0,0}
\usepackage{listings}
\usepackage{xepersian}
\usepackage[noend]{algpseudocode}


%--------------------- page settings ----------------------

\settextfont[Scale=1.1]{XB Niloofar}
\setdigitfont[Scale=1.1]{XB Niloofar}
\defpersianfont\sayeh[Scale=1.1]{XB Niloofar}
\addtolength{\textheight}{3.2cm}
\addtolength{\topmargin}{-22mm}
\addtolength{\textwidth}{3cm}
\addtolength{\oddsidemargin}{-1.5cm}


%------------------------ Environments ------------------------------------

\newtheorem{قضیه}{قضیه}
\newtheorem{لم}{لم}
\newtheorem{مشاهده}{مشاهده}
\newtheorem{تعریف}{تعریف}


%-------------------------- Notations ------------------------------------
\renewcommand{\labelitemi}{$\bullet$}
\newcommand{\IR}{\ensuremath{\mathbb{R}}} 
\newcommand{\IZ}{\ensuremath{\mathbb{Z}}} 
\newcommand{\IN}{\ensuremath{\mathbb{N}}} 
\newcommand{\IS}{\ensuremath{\mathbb{S}}} 
\newcommand{\IC}{\ensuremath{\mathbb{C}}} 
\newcommand{\IB}{\ensuremath{\mathbb{B}}} 

\newcommand{\bR}{\mathbb{R}}
\newcommand{\cB}{\mathcal{B}}
\newcommand{\cO}{\mathcal{O}}
\newcommand{\cG}{\mathcal{G}}
\newcommand{\rM}{\mathrm{M}}
\newcommand{\rC}{\mathrm{C}}
\newcommand{\rV}{\mathrm{V}}

\newcommand{\lee}{\leqslant}
\newcommand{\gee}{\geqslant}
\newcommand{\ceil}[1]{{\left\lceil{#1}\right\rceil}}
\newcommand{\floor}[1]{{\left\lfloor{#1}\right\rfloor}}
\newcommand{\prob}[1]{{\mbox{\tt Pr}[#1]}}
\newcommand{\set}[1]{{\{ #1 \}}}
\newcommand{\seq}[1]{{\left< #1 \right>}}
\newcommand{\provided}{\,|\,}
\newcommand{\poly}{\mbox{\rm poly}}
\newcommand{\polylog}{\mbox{\rm \scriptsize polylog}\,}
\newcommand{\comb}[2] {\left(\!\!\begin{array}{c}{#1}\\{#2}\end{array}\!\!\right)}




\newcounter{probcnt}
\newcommand{\مسئله}[1]{\stepcounter{probcnt}{
 	\bf \arabic{probcnt}$\mbox{\bf{.}}$ \ #1}}

\newcommand{\fqed}[1]{\leavevmode\unskip\nobreak\quad\hspace*{\fill}{\ensuremath{#1}}}

\newenvironment{اثبات}
	{\begin{trivlist}\item[\hskip\labelsep{\em اثبات.}]}
	{\fqed{\square}\end{trivlist}}

\newenvironment{حل}
	{\begin{trivlist}\item[\hskip\labelsep{\bf حل.}]}
	{\fqed{\blacktriangleright}\end{trivlist}}

\ifdefined\hidesols
	\newsavebox{\trashcan} % uncomment the following line to hide solutions
	\renewenvironment{حل}{\begin{latin}\begin{lrbox}{\trashcan}}{\end{lrbox}\end{latin}}
\fi


%------------------------- Header -----------------------------

%------------------------- Header -----------------------------

\newcommand{\سربرگ}[3]{
	\parindent=0em
	
	
	\begin{shaded}
		
		\rightline{ 
			\makebox[8em][c]{
				\includegraphics[height=3cm]{aut.png}
		}} \ \
	\\[-3em] 
	\centerline{ مدرس :دکتر امیر مزلقانی } 
	\\[-6.5em]
	\centerline{\large \bf جبرخطی کاربردی }
	\\[0.05em]
	\centerline {\bf نیمسال دوم 98-97}	
	\\[-6.4em]
		\leftline{ 
		\makebox[8em][c]{
			\includegraphics[height=3.2cm]{ceit.jpg}
	}}
	
	
		

		\hrule height .12em
		
		\normalsize
		\vspace{1mm} #1
		\hfill \small  #3
		\vspace{1mm} 
		\hrule height .1em
		
		\vspace{-0.5em} 
		\hfill {\sayeh\large #2} \hfill
	\end{shaded}
	
{\bf توجه !!!}	
\begin{itemize}
	\item 
	سوالات زیر مربوط به فصل اول درس جبر خطی کاربردی با موضوع ((معادلات خطی در جبر خطی )) می باشد که شامل 5 سوال تئوری و 2 سوال عملی است
	\item 
	سوالات را به دقت و مطالعه و به صورت خوانا و مرتب بنویسید
	\item 
	مهلت ارسال این تمرین ساعت 
	\color{red}{23:55 روز شنبه 97/12/18}
	می باشد. 
\end{itemize}
}





\begin{document}
\سربرگ{مجموعه سوالات فصل 2و3 (جبر ماتریسی و دترمینان)}{}{}


\مسئله{} 
ابتدا دترمینان ماتریس های زیر را بیابید سپس با استفاده از دو روش استفاده از ماتریس های مقدماتی و ماتریس الحاقی (adjugate) وارون آن هارا در صورت وجود بیابید.
 
$$\begin{bmatrix}
1&1&3\\
-2&2&1\\
0&1&1
\end{bmatrix}\quad 
\begin{bmatrix}
-2&-7&-9\\
2&5&6\\
1&3&4\\
\end{bmatrix}\quad 
\begin{bmatrix}
1&-2&1\\
4&-7&3\\
-2&6&-4\\
\end{bmatrix}\quad
\begin{bmatrix} 
0&1&1&\cdots&1\\
1&0&1&\cdots&1\\
1&1&0&\cdots&1\\
\vdots&\vdots&\vdots&\ddots&\vdots\\
1&1&1&\cdots&0
\end{bmatrix}_{n\times n}$$$$
\begin{bmatrix} 
0&a_1&1&\cdots&0\\
0&0&a_2&\cdots&0\\
\vdots&\vdots&\ddots&\ddots&0\\
0&0&0&\cdots&a_{n-1}\\
a_n&0&0&\cdots&0
\end{bmatrix}_{n\times n}
$$
$a_i$
ها پارامتر و غیر صفر هستند.


\مسئله{}
درستی یا نادرستی گزاره های زیر را مشخص کنید برای گزاره های نادرست مثال نقض و برای گزاره های درست اثبات ارائخ دهید
\begin{enumerate}
\item 
اگر 
$A,B$
ماتریس های 
$n\times n $
و 
$AB-BA=A$
آنگاه 
$A$
معکوس پذیر نیست.
\item 
	ماتریسی را پاد متقارن گوییم که برابر قرینه ترانهاده اس باشد ،اگر 
	$A$
	ماتریسی 
	$n\times n$
	باشد آنگاه اگر 
	$n$
	زوج باشد
	$det A=0$.
	\item 
	اگر 
	$A$
	یک ماتریس بالا مثلثی باشد اگر هیچ یک از درایه های روی قطر اصلی آن برابر صفر نباشد آنگاه هم ارز سطری 
	$I_n$
	است.
\end{enumerate}
\مسئله{}
فرض کنید 
$A$
یک ماتریس 
$n\times n$
باشد که درایه های آن 
$-1,1$
تشکیل شده باشد ،ثابت کنید 
$2^{n-1}$
عاد می کند 
$det A$
را .

\مسئله {}
اگر یک ماتریس 
$A$
یک ماتریس 
$n\times n $
باشد 
گزاره های زیر را ثابت کنید:
\begin{enumerate}
\item 
	$adj(A^t)=(adj A)^t$
\item 
$adj  A^{-1}=(adj A)^{-1}$
\item 
اگر 
$A$
ماتریس قطری باشد انگاه 
$adj A$
نیز قطری است.
\item 
نشان دهید اگر 
$I-AB$
معکوس پذیر باشد آنگاه 
$I-BA$
نیز معکوس پذیر است.
\item 
نشان دهید :
$|adj(adj(A))|=|A|^{(n-1)^2}$
	
\end{enumerate}
\مسئله {}
فرض کنید 
$A$
یک ماتریس 
$n\times n$
باشد و 
$A=a_{(ij)}$
به طوری که 
به ازای هر 
$1\leq i\leq n $
داشته باشیم 
\\
$\sum_{j=1}^{n}a_{ij}=a$
که 
$a$
مستقل از درایه های ماتریس است.اگر 
$A^2=I$
مقدار 
$a$
را محاسبه کنید.

\مسئله {}
فرض کنید 
$A,B,C$
ماتریس هایی 
$n\times n$
 باشند ،به ازای چه 
 $n$
 هایی رابطه زیر برقرار است :
 $$(AB-BA)^2C=C(AB-BA)^2$$
 \مسئله {}
 دترمینان ماتریس های زیر را بیابید:
 $$\begin{bmatrix}
 a+b&ab&0&\cdots&0&0\\
 1&a+b&ab&\cdots&0&0\\
 0&1&a+b&\cdots&0&0\\
 \cdots&\cdots&\cdots&\cdots&\cdots&\cdots\\
 0&0&0&\cdots&a+b&ab\\
 0&0&0&\cdots&1&a+b
 \end{bmatrix}_{n\times n}\quad
 \begin{bmatrix}
 1&1&1&\cdots&1\\
 a_1&a_2&a_3&\cdots&a_n\\
 a_1^2&a_2^2&a_3^2&\cdots&a_n^2\\
 \cdots&\cdots&\cdots&\cdots&\cdots\\
 a_1^{n-1}&a_2^{n-1}&a_3^{n-1}&\cdots&a_n^{n-1}
\end{bmatrix}_{n\times n}\\
 $$
 $$
 \begin{bmatrix}
 1&1&1&\cdots&1\\
 1&2&2&\cdots&2\\
 1&2&3&\cdots&3\\
 \vdots&\vdots&\vdots&\ddots&\vdots\\
 1&2&3&\cdots&n
 \end{bmatrix}_{n\times n}\quad
  $$
 \مسئله {}فرض کنید تمامی ماتریس های زیر وارون پذیر باشند نشان دهید :
 $$(A-B)^{-1}=A^{-1}+A^{-1}(B^{-1}-A^{-1})^{-1}A^{-1}$$
 همچنین در حالت خاص نشان دهید :
 $$(I+A)^{-1}=I-(A^{-1}+I)^{-1}$$
 و نشان دهید:
 $$|(I+A)^{-1}+(I+A^{-1})^{-1}|=1$$
 \مسئله{}
جواب دستگاه معادلات زیر را به روش کرامر بیابید:
\begin{equation*}
\left\{
\begin{array}{rl}
x_1+x_2&=3\\
-3x_1+2x_3&=0\\
x_2-2x_3&=2
\end{array} \right.\qquad
\left\{
\begin{array}{rl}
x_1+3x_2+x_3&=4\\
-x_1+2x_3&=2\\
3x_1+x_2&=2
\end{array} \right.
\end{equation*}
\مسئله {}دستگاه معادلات  زیر را با استفاده از تجزیه 
$LU$
حل کنید همچینین وارون ماتریس افزوده را با تجزیه 
$LU$
به دست آورید.
\begin{equation*}
\left\{
\begin{array}{rl}
	x_1+3x_2+3x_4&=8\\
	2x_1+X_2-2x_3+x_4&=7\\
	3x_1-x_2+x_3+2x_4&=14\\
	-x_1+2x_2+3x_3-x_4&=-7
\end{array} \right.
\end{equation*}
\مسئله{}
فرض کنید 
$T:\mathbb{R}^3\longrightarrow \mathbb{R}^3$
یک تبدیل خطی با ماتریس 
$\begin{bmatrix}
a&0&0\\
0&b&0\\
0&0&c
\end{bmatrix}$
باشد که 
$a,b,c$
مقادیری مثبت باشند،
$S$
را کره واحد در نظر بگیرید که سطح آن با معادله 
$x_1^2+x_2^2+x_3^2=1$
محدود شده است،
\begin{enumerate}
	\item 
نشان دهید 
$T(S)$
با بیضی به معادله 
$\frac{x_1^2}{a^2}+\frac{x_2^2}{b^2}+\frac{x_3^2}{c^2}=1$
محدود شده است.
\item 
با این فرض که حجم کره واحد 
$4\pi/3$
است حجم بیضی مطرح شده در قسمت 
$1$
را بیابید.
\end{enumerate}
\end{document}